\begin{abstract}
In modern industrial systems, acquiring labeled training data for all potential fault types remains prohibitively expensive and often impractical, as novel or rare faults may never have been observed during operation. Zero-shot learning (ZSL) offers a promising paradigm to address this challenge by enabling fault classifiers to recognize unseen fault categories through semantic knowledge transfer from seen classes. However, existing ZSL-based fault diagnosis methods, particularly those relying on Generative Adversarial Networks (GANs), suffer from training instability, mode collapse, and significant domain shift between generated and real features. To overcome these limitations, this paper proposes CESM-Diff (CLIP-Enhanced Semantic Manifold Diffusion), a novel framework that leverages pretrained CLIP text encoders and a dual-path denoising diffusion probabilistic model to synthesize highly discriminative features for unseen fault classes. CESM-Diff introduces three key innovations: (1) a Semantic Manifold Interpolation (SMI) module that constructs continuous fault representations in the CLIP embedding space, enabling smooth transitions between fault severity levels; (2) a dual-path diffusion architecture that simultaneously refines feature-level and semantic-level representations through cross-path attention mechanisms, ensuring semantic-geometric consistency; and (3) an Attribute Consistency Loss ($\mathcal{L}_{acl}$) that constrains generated features to remain physically plausible and semantically aligned with their target fault descriptions. Extensive experiments on three widely-used industrial benchmarks---the Tennessee Eastman Process (TEP), Hydraulic System, and Case Western Reserve University (CWRU) Bearing datasets---demonstrate that CESM-Diff achieves state-of-the-art zero-shot fault diagnosis accuracy of 91.3\% on CWRU, 82.4\% on TEP, and 79.8\% on Hydraulic, significantly outperforming existing GAN-based and diffusion-based methods. Comprehensive ablation studies and t-SNE visualizations further validate the effectiveness of each proposed component.
\end{abstract}
